\documentclass[11pt,a4paper]{article}
\usepackage[T1]{fontenc}
\usepackage[utf8]{inputenc}
\usepackage[polish]{babel}
\usepackage{lmodern}
\usepackage{geometry}
\usepackage{microtype}
\usepackage{booktabs}
\usepackage{amsmath,amssymb}
\usepackage[hidelinks]{hyperref}
\geometry{margin=26mm}
\setlength{\parskip}{0.55em}
\setlength{\parindent}{0pt}
\title{Nasilenie objawów lękowych a elastyczność poznawcza i strategie podejmowania decyzji w warunkach niepewności}
\author{Mateusz Pomianek}
\date{Projekt badania, wersja 0.4.0 \\ 26 listopada 2026}
\begin{document}
\maketitle
\begin{center}
\textit{Anxiety Symptom Severity, Cognitive Flexibility, and Decision-Making Strategies Under Uncertainty}
\end{center}
\begin{abstract}
Celem projektowanego badania jest ocena związku między nasileniem objawów lękowych, kosztem przełączania reguł zadaniowych i strategiami wyboru w warunkach niepewności. W badaniu przekrojowym z komponentem behawioralnym proponuje się połączenie zwalidowanego kwestionariusza nasilenia objawów lękowych, zadania task-switching i niezależnego zadania decyzji probabilistycznych. Głównym estymandem jest zmiana kosztu przełączania przypadająca na pięciopunktową różnicę nasilenia objawów przy zadanym zestawie współzmiennych. Oddzielnie zostanie oszacowany związek z preferencją ryzyka lub niejednoznaczności, o ile parametry zadania umożliwią ich identyfikację.

\textbf{Status:} prospektywny szkic naukowy (\texttt{IDEA\_REVIEW}); nie przeprowadzono badania, nie uzyskano wyników ani zatwierdzenia etycznego.
\end{abstract}
\tableofcontents
\section{Wprowadzenie i cel}
Celem projektu jest sprawdzenie, czy nasilenie objawów lękowych współwystępuje z kosztami przełączania reguł i strategiami wyboru przy ryzyku i niejednoznaczności. Badanie zakłada rozdzielenie kontroli wykonawczej, trafności wykonania, ostrożności reakcji i preferencji decyzyjnych. Wszystkie analizy mają charakter prospektywny; plan obserwacyjny nie pozwala wnioskować o przyczynowości.

\section{Ramy teoretyczne i model zależności}
\subsection{Teoria kontroli uwagi}
Teoria kontroli uwagi zakłada, że nasilony lęk może zmniejszać efektywność kontroli wykonawczej, choć poprawność reakcji może pozostać zachowana dzięki dodatkowej mobilizacji zasobów \cite{Eysenck2007}. Zatem spowolnienie reakcji i wzrost odsetka błędów należy traktować jako odrębne wyniki. Konieczne jest także rozróżnienie zmiany reguły zadaniowej, szybkości bazowej oraz zachowania ostrożnego.

\subsection{Elastyczność poznawcza}
Zdolność przełączania jest powiązana z innymi funkcjami wykonawczymi, ale nie jest z nimi tożsama \cite{Miyake2000}. Koszt przełączania zależy od wskazówek, przygotowania do zadania, kolejności prób i powtarzania bodźców \cite{Monsell2003}. Pojedyncze zadanie jest wskaźnikiem behawioralnym, a nie bezpośrednią miarą neuronalnej elastyczności.

\subsection{Ryzyko a niejednoznaczność}
Ryzyko zakłada znajomość prawdopodobieństw, niejednoznaczność zaś ich niepełną znajomość. Nietolerancja niepewności, rozważana w badaniach nad lękiem, nie jest tym samym co niechęć do wariancji wypłat \cite{Carleton2016}. Z tego powodu należy konstruować odrębne kontrasty decyzyjne.


\subsection{Stan dowodów i luka badawcza}
Metaanaliza badań nad kontrolą uwagi wskazuje na znaczenie komponentów hamowania i przełączania oraz rozróżnienia efektywności i trafności \cite{Shi2019}. Metaanaliza dotycząca osób z rozpoznanym GAD wskazuje również na możliwość obniżonej trafności w zadaniach elastyczności poznawczej \cite{GAD2025}. Wyników klinicznych nie należy bezpośrednio przenosić na zależności wymiarowe w populacji ogólnej. Wong i Yu \cite{Wong2026} nie znaleźli związku kosztów przełączania afektywnego z lękiem po korekcie wielokrotnych porównań. Uzasadnia to testowanie hipotez kierunkowych przy dwustronnym wnioskowaniu i dopuszczenie wyniku zerowego.

Przegląd Hartley i Phelps \cite{Hartley2012} łączy lęk z procesami wyceny i przewidywania zagrożenia, lecz nie uzasadnia utożsamiania wszystkich form awersji do niepewności. Badania nad wartością subiektywną wskazują na częściowo wspólne mechanizmy wyceny ryzyka i niejednoznaczności \cite{Levy2010}. Niniejszy projekt integruje dwa paradygmaty w tej samej próbie i testuje odrębne kontrasty, bez przypisywania zachowaniu swoistych mechanizmów neuronalnych.

Istotnym ograniczeniem jest paradoks rzetelności: stabilny efekt eksperymentalny na poziomie grupy może mieć małą rzetelność różnic indywidualnych \cite{Hedge2018}. Dlatego pilotaż musi obejmować nie tylko średni koszt przełączania, ale także jego rzetelność i zmienność międzyosobniczą. Rejestracja czasu reakcji online wymaga kontroli przeglądarki, zdarzeń utraty fokusu i stabilności pomiaru \cite{Bridges2020}.
\subsection{Operacjonalizacja konstruktu lęku}
GAD-7 mierzy nasilenie objawów w ostatnich dwóch tygodniach; nie jest miarą chwilowego lęku ani samodzielnym rozpoznaniem klinicznym \cite{Spitzer2006}. Dostępne są dane walidacyjne polskiej wersji w populacji pracującej i populacji ogólnej \cite{GADPL2023,GADPL2025}. Wynik ciągły 0--21 będzie analizowany bez arbitralnej dychotomizacji. Nasilenie objawów depresyjnych zostanie uwzględnione w analizie warunkowej; ze względu na współwystępowanie objawów należy również pokazać oszacowanie bez tej współzmiennej.

\section{Pytania i planowane hipotezy szczegółowe}
\begin{description}
\item[H1, główna:] Większy wynik GAD-7 wiąże się z większym kosztem przełączania po kontroli szybkości bazowej i uprzednio wybranych współzmiennych.
\item[H2, drugorzędowa:] Warunkowo względem różnicy wartości oczekiwanej i wariancji wyższy GAD-7 wiąże się z rzadszym wyborem opcji o większej wariancji.
\item[H3, drugorzędowa:] Warunkowo względem nominalnego kontrastu wypłat i szerokości przedziału wyższy GAD-7 wiąże się z rzadszym wyborem opcji o nieznanych prawdopodobieństwach.
\item[H4, eksploracyjna:] Związki z decyzjami pozostają po uwzględnieniu indywidualnego kosztu przełączania.
\end{description}
Hipotezy nie stanowią stwierdzenia o wynikach. Kierunek, definicje i analiza wielokrotności muszą zostać zamrożone przed rekrutacją.

\section{Szczegółowa procedura eksperymentalna}
\subsection{Próba i warunki techniczne}
Planowane jest badanie online osób pełnoletnich. Zadanie szybkościowe powinno być wykonywane na komputerze z klawiaturą. Uczestnik otrzyma instrukcję dobrowolności i możliwości rezygnacji. Zostaną zebrane wyłącznie minimalne dane o wieku, śnie, porze wykonania zadania i konfiguracji urządzenia. Kryteria wykluczeń technicznych i liczebność próby wymagają osobnego uzasadnienia.

\subsection{Sekwencja}
Kolejność obejmie zgodę, krótkie pytania kontrolne, GAD-7, trening i właściwe próby przełączania, zadanie decyzji oraz debriefing. Kolejność dwóch zadań zostanie losowo zrównoważona. Wstępna liczba bloków, prób i czasy są określone poniżej; wartości ostateczne zostaną ustalone w pilotażu, bez analizowania hipotez na próbie konfirmacyjnej.

\subsection{Paradygmat przełączania}
Uczestnik będzie klasyfikować bodźce według dwóch reguł, np. parzystości i wielkości liczby. Wskazówka zadaniowa określi aktualną regułę. Należy zrównoważyć typy prób switch i repeat, rodzaje bodźców oraz zgodność odpowiedzi. Miarą podstawową będzie różnica median czasu poprawnych reakcji:
\begin{equation}
SC_i=\operatorname{median}(RT^{\mathrm{correct}}_{i,\mathrm{switch}})
-\operatorname{median}(RT^{\mathrm{correct}}_{i,\mathrm{repeat}}).
\end{equation}
Uzupełniająco należy raportować różnice w błędach, względny koszt RT i jakość danych na poziomie prób. Selekcja wyłącznie reakcji poprawnych wymaga osobnej analizy wrażliwości.

\subsection{Zadanie decyzyjne}
W wariancie ryzyka obie opcje mają jawne prawdopodobieństwa oraz wypłaty. Wartość oczekiwana wynosi:
\begin{equation}
EV_j=\sum_{k=1}^{K_j}p_{jk}x_{jk}.
\end{equation}
Dla H2 należy kontrolować wartość oczekiwaną, zmieniając wariancję. W wariancie niejednoznaczności wybrane prawdopodobieństwa pozostają nieznane; uczestnik musi to rozumieć. Pilotaż ustali, czy preferencję niejednoznaczności da się oddzielić od atrakcyjności wypłaty.

\section{Główny estymand i modele statystyczne}
Główny estymand to średnia różnica kosztu przełączania w milisekundach związana z pięciopunktową różnicą GAD-7, warunkowo względem ustalonych współzmiennych. Jednostka pięciu punktów ma charakter opisowy, a nie diagnostyczny.
\begin{equation}
SC_i=\beta_0+\beta_1\frac{GAD7_i}{5}
+\boldsymbol{\gamma}^{\mathsf T}\mathbf Z_i+\varepsilon_i.
\end{equation}
Współzmienne będą ograniczone do uzasadnionych czynników, wieku, czasu snu, nasilenia objawów depresyjnych, typu urządzenia, kolejności zadań i bazowego RT z odrębnych bloków jednorodnych. Raportowane będą oszacowania, przedziały ufności oraz diagnostyka reszt. Regresja odporna na obserwacje odstające będzie analizą wrażliwości.

Dodatkowo model mieszany na poziomie prób uwzględni interakcję przełączenia z GAD-7 i intercept uczestnika. Odrębny model logistyczny oceni poprawność reakcji, aby wykryć kompromis szybkość--trafność. Wybory decyzyjne można analizować modelem logistycznym z efektem uczestnika, różnicą wartości oczekiwanej, wariancją i wskaźnikiem niejednoznaczności. Parametry nieidentyfikowalne nie będą interpretowane.

H1 stanowi pojedynczy test główny. Dla H2 i H3 przewidziano korektę Holma, a H4 pozostanie eksploracyjna. Reguły obsługi braków danych i odrzucania prób zostaną ustalone przed analizą.

\section{Pilotaż, wielkość próby i kontrola jakości}
Liczebność próby zostanie uzasadniona przy użyciu minimalnego efektu istotnego naukowo oraz symulacji mocy lub precyzji estymacji. Pilotaż posłuży do oszacowania wariancji kosztu przełączania, odsetka prób nieważnych, rozkładu GAD-7, zrozumiałości instrukcji i identyfikowalności wyborów. Dane używane do dostrajania paradygmatu nie będą automatycznie łączone z próbą konfirmacyjną.

System badawczy powinien zapisywać identyfikator wersji bodźca, regułę, wskazówkę, blok, poprawność, czas reakcji i parametry wyboru. Analizy wrażliwości obejmą przeglądarkę, urządzenie, kolejność zadań, trafność, przekształcenia RT i ewentualne odstępstwa od planu. Publicznie udostępniane będą co najwyżej kod, słownik zmiennych i dane syntetyczne.


\section{Projekt operacyjny (do zatwierdzenia po pilotażu)}
Badanie przekrojowe, zdalne, polskojęzyczne, z rekrutacją osób od 18. roku życia. Niezależnie od wyniku GAD-7 uczestnicy wykonują te same zadania. Nie będzie selekcji na podstawie samoopisu rozpoznania. Przewidywany czas sesji to 30--40 minut, do sprawdzenia w pilotażu. Kolejność zadań i przypisania odpowiedzi do klawiszy będą równoważone; ziarno losowania będzie zapisywane. Nie przewiduje się indukcji lęku ani wypłat pieniężnych.

\subsection{Parametry zadania przełączania}
Dwie reguły dotyczą liczb 1--9 z wyłączeniem 5: parzystość i relacja względem 5. Dwa różne sygnały na każdą regułę pozwolą wymusić zmianę sygnału także w próbach repeat; odpowiedzi lewa/prawa będą równoważone. Proponuje się dwa bloki jednorodne po 24 próby do niezależnego pomiaru bazowego RT oraz cztery bloki mieszane po 48 prób. Pierwsza próba każdego bloku mieszanego nie ma zdefiniowanego przejścia i nie wchodzi do estymacji kosztu. Łącznie 188 potencjalnych przejść, po 94 switch i repeat. Sekwencje mają zrównoważyć regułę, zgodność odpowiedzi, powtórzenie bodźca i typ przejścia, z udokumentowaną tolerancją odchyleń. Wstępne parametry: wskazówka 300 ms, okno odpowiedzi 2000 ms, przerwa 500 ms. Trening: 24 próby, minimum 75\% poprawności, jedno powtórzenie. Progi i czasy podlegają zamrożeniu po pilotażu.

\subsection{Parametry zadania wyboru}
Wstępnie 48 prób ryzyka i 48 prób niejednoznaczności, z losową stroną opcji, równoważeniem kolejności domen i próbami sprawdzającymi zrozumienie. W próbach ryzyka przykładem porównania jest loteria o wypłatach 4 i 8 punktów oraz loteria o wypłatach 0 i 12 punktów, w obu przypadkach przy prawdopodobieństwie 0,5. Obie mają EV = 6, wariancje odpowiednio 4 i 36. Pozostałe poziomy EV będą zmieniane niezależnie od wariancji, aby zidentyfikować ich efekty. W próbach niejednoznaczności porównywane są loterie o kontrolowanych wypłatach (różnica nominalna $-2,0,2$ punktów): jawne prawdopodobieństwo 0,5 kontra przedział symetryczny wokół 0,5 (np. 0,3--0,7). Środek przedziału jest tylko nominalną wartością do konstrukcji bodźca, nie znanym prawdopodobieństwem ani rzeczywistym EV. Szerokość przedziału będzie zmieniana. Wybory hipotetyczne nie są równoznaczne z wyborami motywowanymi ekonomicznie; wnioski ograniczymy do tego paradygmatu.
\begin{equation}
V_j=\sum_k p_{jk}(x_{jk}-EV_j)^2
\end{equation}
W obu zadaniach wyboru zmienną zależną jest wybór zdefiniowanej opcji docelowej (bardziej zmienna lub niejednoznaczna), nie wybór strony ekranu.

\subsection{Plan analiz i kontrola jakości}
Analiza H1: OLS z odpornymi błędami standardowymi HC3, dwustronne $\alpha=0{,}05$ i 95\% przedział ufności. Jednostka obserwacji to uczestnik. Współzmienne: wiek, godziny snu, wynik depresyjny, system operacyjny, rodzina przeglądarki, kolejność i niezależne bazowe RT; rzadkie kategorie zostaną połączone według reguły zamrożonej w pilotażu. Bazowe RT nie będzie wyliczane z prób repeat użytych do kosztu, aby ograniczyć mechaniczne sprzężenie wyniku ze współzmienną. Analiza wrażliwości obejmie model bez bazowego RT oraz bez wyniku depresyjnego, alternatywną transformację logarytmiczną RT i model prób z losowym interceptem i, jeżeli identyfikowalny, nachyleniem switch dla uczestnika.
\begin{equation}
\operatorname{logit}P(Y_{it}=1)=a+u_i+bGAD7_i/5+c\Delta EV_{it}+d\Delta V_{it}+\boldsymbol{\gamma}^{\mathsf T}\mathbf Z_i
\end{equation}
Zakłada się $u_i\sim\mathcal N(0,\sigma_u^2)$, niezależnie od współzmiennych. Współczynnik $b$ jest wspólnym warunkowym logarytmem ilorazu szans na 5 punktów GAD-7 przy ustalonych kontrastach, także dla kontrastu wypłat równego zero. Interakcje GAD-7 z kontrastami są eksploracyjne; współczynnik nie identyfikuje strukturalnego parametru awersji. Powyższy model dotyczy prób ryzyka. W modelu niejednoznaczności zamiast $\Delta V$ i rzeczywistego $\Delta EV$ uwzględnia się nominalny kontrast wypłaty i szerokość przedziału prawdopodobieństw. Dla obu modeli przewidziano współzmienne z H1 z wyjątkiem bazowego RT, dodawanego eksploracyjnie; dodatkowo kolejność domen decyzji. H2 i H3: dwustronne testy współczynników GAD-7 z korektą Holma; raportowanie ilorazów szans i marginalnych prawdopodobieństw. H4 nie będzie traktowana jako analiza mediacji przyczynowej.

Wstępne kryteria: wykluczenie prób z utratą fokusu, błędem technicznym, RT poniżej 150 ms lub przekroczeniem 2000 ms; błędy merytoryczne pozostają w analizie trafności. Do H1 wymagane minimum 75\% poprawności oraz po 50 ważnych poprawnych prób switch i repeat oraz 24 poprawne bazowe. Timeout liczy się jako niepoprawna odpowiedź w trafności. Próby po błędzie pozostają w analizie głównej, a ich wyłączenie stanowi analizę wrażliwości. Nie stosować progów zależnych od GAD-7 ani kierunku efektu. Brak któregokolwiek elementu GAD-7 oznacza brak pełnego wyniku; brak współzmiennych będzie raportowany, z analizą kompletnych przypadków jako główną i wielokrotną imputacją współzmiennych jako analizą wrażliwości, jeśli uzasadniona. Odrzucenie zadania nie wyklucza automatycznie drugiego zadania. Raportować przepływ uczestników i rzetelność split-half kosztu z korektą Spearmana--Browna oraz przedziałem bootstrap. W miarę możliwości oszacować powtarzalność w osobnej sesji pilotażu.

\subsection{Uzasadnienie liczebności}
Wstępnie przyjęto minimalny efekt częściowy $|r|=0{,}15$, moc 80\% i dwustronne $\alpha=0{,}05$. Przybliżenie Fishera dla korelacji częściowej z $k$ współzmiennymi:
\begin{equation}
n\simeq k+3+\left(\frac{z_{1-\alpha/2}+z_{0{,}80}}{\operatorname{atanh}(r)}\right)^2.
\end{equation}
Dla ilustracyjnego $k=8$ potrzeba około 355 kompletnych obserwacji; przy 20\% utracie plan rekrutacji wynosi 444 osoby. Dla $r=0{,}20$ potrzeba około 202 kompletnych obserwacji; dla $r=0{,}10$ około 791. Jest to planistyczne przybliżenie, a nie oszacowanie efektu z danych. Ostateczne $k$ oznacza liczbę stopni swobody współzmiennych (w tym kodowanie kategorii). Pilotaż 30--50 osób niezależnych od próby głównej posłuży do symulacji mocy z założoną wariancją RT, rozkładem GAD-7, brakami i rzetelnością; efekty H2/H3 wymagają odrębnej symulacji. Nie zmieniać założeń na podstawie zaobserwowanego efektu GAD-7 w pilotażu. W razie braku możliwości rekrutacji odpowiedniej próby ograniczyć cel do badania estymacyjnego z jawnie określoną precyzją.

\subsection{Etyka i zarządzanie danymi}
Przed rekrutacją wymagane są decyzja właściwej komisji, uzgodnienie administratora danych i podstaw przetwarzania danych o zdrowiu, zatwierdzone polskie wersje skal oraz informacja dla uczestnika. Planowana skala depresji PHQ-9 wymaga zatwierdzonej procedury reagowania na pozycję dotyczącą samouszkodzenia; alternatywa bez tej pozycji wymaga odrębnej weryfikacji trafności i licencji oraz zmiany protokołu przed zamrożeniem. Formularz nie może sugerować monitorowania w czasie rzeczywistym, diagnozy ani udzielania pomocy klinicznej. Uczestnik może przerwać badanie; debriefing zawiera informacje o dostępnej pomocy. Dane pseudonimizowane, bez nazwisk, IP i swobodnych opisów zdrowotnych, z ograniczeniem metadanych technicznych do potrzeb QA. Dane i klucz identyfikacyjny (jeżeli potrzebny) poza publicznym repozytorium, oddzielnie, w zatwierdzonej infrastrukturze. Ustalić okres retencji, zasady wycofania i zakres udostępniania przed uruchomieniem.

\section{Scenariusze interpretacji}
Jeżeli koszt przełączania będzie większy przy nasilonych objawach, wynik będzie zgodny z hipotezą o obniżonej efektywności wykonywania zadania, ale nie zidentyfikuje mechanizmu neuronalnego ani zależności przyczynowej. Brak związku może odzwierciedlać rzeczywiście mały efekt, zbyt małą precyzję albo niewystarczającą trafność paradygmatu. Wzorce wyborów należy interpretować osobno dla ryzyka i niejednoznaczności; nie stanowią one samodzielnej diagnozy klinicznej.

\section{Wyniki}
Na dzień 26 listopada 2026 r. nie przedstawiono danych empirycznych. Po przeprowadzeniu badania należy zamieścić pełne wyniki, w tym wyniki niezgodne z hipotezami, dane o odrzuceniach, niepewność oszacowań oraz analizy wrażliwości.
\section{Warunki przejścia do badania}
\begin{enumerate}
\item Zweryfikować istniejącą literaturę i przyporządkować cytowania do twierdzeń.
\item Ustalić finalne narzędzia, wersje językowe, licencje i procedurę pilotażową.
\item Zatwierdzić estymand, hipotezy, liczebność próby oraz plan kontroli błędów wielokrotności.
\item Rozstrzygnąć wymogi etyczne i ochronę danych, a następnie przygotować prerejestrację.
\item Dopiero wtedy rozpocząć rekrutację i gromadzenie materiału badawczego.
\end{enumerate}
\section*{Bibliografia}
Poniżej zestawiono literaturę założycielską i uzupełniający przegląd ukierunkowany. 

\begin{thebibliography}{99}
\bibitem{Eysenck2007} Eysenck, M. W., Derakshan, N., Santos, R., Calvo, M. G. (2007). Anxiety and cognitive performance: Attentional control theory. \textit{Emotion}, 7(2), 336--353. \href{https://doi.org/10.1037/1528-3542.7.2.336}{doi:10.1037/1528-3542.7.2.336}.
\bibitem{Miyake2000} Miyake, A., Friedman, N. P., Emerson, M. J., Witzki, A. H., Howerter, A., Wager, T. D. (2000). The unity and diversity of executive functions and their contributions to complex ``Frontal Lobe'' tasks: A latent variable analysis. \textit{Cognitive Psychology}, 41(1), 49--100. \href{https://doi.org/10.1006/cogp.1999.0734}{doi:10.1006/cogp.1999.0734}.
\bibitem{Monsell2003} Monsell, S. (2003). Task switching. \textit{Trends in Cognitive Sciences}, 7(3), 134--140. \href{https://doi.org/10.1016/S1364-6613(03)00028-7}{doi:10.1016/S1364-6613(03)00028-7}.
\bibitem{Spitzer2006} Spitzer, R. L., Kroenke, K., Williams, J. B. W., Löwe, B. (2006). A brief measure for assessing generalized anxiety disorder: The GAD-7. \textit{Archives of Internal Medicine}, 166(10), 1092--1097. \href{https://doi.org/10.1001/archinte.166.10.1092}{doi:10.1001/archinte.166.10.1092}.
\bibitem{Carleton2016} Carleton, R. N. (2016). Fear of the unknown: One fear to rule them all? \textit{Journal of Anxiety Disorders}, 41, 5--21. \href{https://doi.org/10.1016/j.janxdis.2016.03.011}{doi:10.1016/j.janxdis.2016.03.011}.

\bibitem{Shi2019} Shi, R., Sharpe, L., Abbott, M. (2019). A meta-analysis of the relationship between anxiety and attentional control. \textit{Clinical Psychology Review}, 72, 101754. \href{https://doi.org/10.1016/j.cpr.2019.101754}{DOI}.
\bibitem{GAD2025} Nguyen, L., Walters, J., Hutchinson, E., Liu, Y., Li, X., Gudmundsson, C. (2025). Executive functioning in individuals with generalized anxiety disorder: A systematic review and meta-analysis. \textit{Journal of Affective Disorders}, 389, 119683. \href{https://doi.org/10.1016/j.jad.2025.119683}{DOI}.
\bibitem{Wong2026} Wong, Y.-S., Yu, J. (2026). Anxiety facilitates disengagement when switching from negative-valenced stimuli: Mediating role of attentional control. \textit{Cognitive, Affective, \& Behavioral Neuroscience}, 26, 104--119. \href{https://doi.org/10.3758/s13415-025-01349-2}{DOI}.
\bibitem{Hartley2012} Hartley, C. A., Phelps, E. A. (2012). Anxiety and decision-making. \textit{Biological Psychiatry}, 72, 113--118. \href{https://doi.org/10.1016/j.biopsych.2011.12.027}{DOI}.
\bibitem{Levy2010} Levy, I., Snell, J., Nelson, A. M., Rustichini, A., Glimcher, P. W. (2010). Neural representation of subjective value under risk and ambiguity. \textit{Journal of Neurophysiology}, 103, 1036--1047. \href{https://doi.org/10.1152/jn.00853.2009}{DOI}.
\bibitem{Hedge2018} Hedge, C., Powell, G., Sumner, P. (2018). The reliability paradox: Why robust cognitive tasks do not produce reliable individual differences. \textit{Behavior Research Methods}, 50, 1166--1186. \href{https://doi.org/10.3758/s13428-017-0935-1}{DOI}.
\bibitem{Bridges2020} Bridges, D., Pitiot, A., MacAskill, M. R., Peirce, J. W. (2020). The timing mega-study: Comparing a range of experiment generators, both lab-based and online. \textit{PeerJ}, 8, e9414. \href{https://doi.org/10.7717/peerj.9414}{DOI}.
\bibitem{GADPL2023} Basińska, B. A., Kwissa-Gajewska, Z. (2023). Psychometric properties of the Polish version of the Generalized Anxiety Disorder Scale (GAD-7) in a non-clinical sample of employees during pandemic crisis. \textit{International Journal of Occupational Medicine and Environmental Health}, 36(4), 493--504. \href{https://doi.org/10.13075/ijomeh.1896.02104}{DOI}.
\bibitem{GADPL2025} Gorgol-Waleriańczyk, J. (2025). Psychometric validation and frequencies of positive screenings in the Polish versions of the Patient Health Questionnaire-9 and Generalized Anxiety Disorder-7 in the general Polish population: a three-sample study. \textit{Current Psychology}, 44, 6495--6503. \href{https://doi.org/10.1007/s12144-025-07646-w}{DOI}.
\end{thebibliography}
\end{document}
